% Generated draft: requires mathematical and source review.
\documentclass[11pt]{article}
\usepackage[margin=25mm]{geometry}
\usepackage{fontspec}
\setmainfont{Latin Modern Roman}
\usepackage{amsmath,amssymb,mathtools,unicode-math}
\setmathfont{Latin Modern Math}
\usepackage{enumitem,needspace,xurl,hyperref}
\setlength{\parindent}{0pt}
\setlength{\parskip}{5pt}
\setlength{\emergencystretch}{3em}
\newcommand{\N}{\mathbb{N}}
\newcommand{\Z}{\mathbb{Z}}
\newcommand{\Q}{\mathbb{Q}}
\newcommand{\R}{\mathbb{R}}
\newcommand{\C}{\mathbb{C}}
\newcommand{\F}{\mathbb{F}}
\newcommand{\A}{\mathbb{A}}
\newcommand{\PP}{\mathbb{P}}
\newcommand{\E}{\mathbb{E}}
\newcommand{\eps}{\varepsilon}
\newcommand{\dd}{\,\mathrm{d}}
\begin{document}
\section*{UnsolvedMath ID: 30001763}
\textbf{Source problem code:} catalogue entry
\subsection{Definitions and mathematical statement}
Let \(Z\) denote the toric degeneration parameter space mentioned in the supplied record, let \(D\subseteq Z\) denote its discriminant, and write
\[
U=Z\setminus D.
\]
The record does not specify the degeneration, the admissible parameter spaces, the discriminant construction, or the topology on \(U\). Its background describes the associated category as a Fukaya category, but does not specify that category or its dependence on the degeneration. Denote this unspecified category by \(\mathcal C\).

The two supplied questions are:
\begin{enumerate}
\item Does the rational homotopy type of \(U\) determine the gaps in the category spectrum of \(\mathcal C\)?
\item Do the ``big loops'' in \(\pi_1(U)\) determine the wall-crossings of \(\mathcal C\)?
\end{enumerate}
Here a based fundamental group would be written \(\pi_1(U,u_0)\) for a chosen point \(u_0\in U\); the record supplies neither a basepoint nor a connectedness assumption.

These questions retain unresolved source terminology. The record does not define the category spectrum, its gaps, the distinguished loops called ``big,'' or the wall-crossings under consideration. It also does not specify what information a loop is to determine, what equivalence of wall-crossings is intended, or what correspondence compares categories belonging to different degeneration data. In the first question, ``determine'' suggests invariance of gap data under an appropriate equivalence of rational homotopy types, but neither the admissible comparisons nor the gap data are defined. Consequently, the supplied evidence supports these two qualified questions, but does not specify a precise universally quantified conjecture.
\subsection{Short English statement}
The source asks whether the rational homotopy type of a toric degeneration parameter space away from its discriminant determines categorical spectral gaps, and whether certain ``big loops'' determine categorical wall-crossings. The background identifies the intended setting as Fukaya categories. The category, spectral invariant, distinguished loops, and wall-crossing data require clarification before these questions become precise conjectures.
\subsection{Sources}
\begin{itemize}
\item Local UnsolvedMath registry: supplied record with permanent ID \(30001763\), exact source problem code \texttt{catalogue entry}, title \emph{Categorical Gaps and Wall-Crossings from Toric Degenerations}; collection: Oberwolfach Reports Open Problems. The proposed year is recorded as 2011; no individual proposer is supplied.
\item Supplied report citation: \emph{Real Enumerative Questions in Complex and Tropical Geometry} (2011), Oberwolfach Reports, pp.\ 1117--1175, DOI \texttt{10.4171/owr/2011/20}. Supplied source URL and literature-source link: \url{https://doi.org/10.4171/owr/2011/20}. No problem-specific page or passage is supplied.
\item Supplied literature assessment, checked \textbf{2026-08-21}: the record labels the problem open and reports that the two Fukaya-category conjectures concerning spectral gaps and wall-crossings are not established in the quoted generality, with no authoritative proof or counterexample found in that curation. This is supplied dated curation, not an independently verified literature assessment. The supplied background expressly requests independent verification of time-sensitive status and cited claims.
\item The registry marks the statement \texttt{corrected\_verified} and reports a manual semantic, typographic, and LaTeX review of its cleanup. This is supplied editorial metadata; it does not resolve the formulation issues identified here or establish mathematical correctness.
\end{itemize}
\end{document}
